% The narrative deck. Built with `make campaigns` -> campaigns.pdf
%
% A fourth document, alongside main.tex (the paper), notebook.tex (the AME
% programme record) and library_notebook.tex (the library arc).
%
% This one exists for a different job from all three: choosing what the paper
% argues, before it argues it. Every campaign in the project measures the same
% two quantities -- how much of a library folds, and how much sequence space it
% covers -- and each one puts a different dial on the trade between them. Here
% each campaign gets one panel, the tree it came from, and the one thing it
% cannot say. The arcs at the end are then picked from the panels rather than
% assembled around a conclusion already chosen.
%
% Figures: paper/figures/campaign_panels.py -> figures/svg/panel_c*.svg
\documentclass[11pt]{article}

% --- Packages ---
\usepackage[utf8]{inputenc}
\usepackage[T1]{fontenc}
\usepackage{lmodern}
\usepackage[margin=1in]{geometry}
\usepackage{microtype}

\usepackage{graphicx}
\usepackage{xcolor}
\usepackage{amsmath, amssymb}
\usepackage{siunitx}
% siunitx drops a leading "+" unless told otherwise, so \SI{+23.0}{points} set
% "23.0 points" -- turning a signed margin into a level and silently reversing
% the reading of any negative-looking cell beside it. There is no unit called
% "points" either; declare it rather than letting siunitx guess.
\sisetup{retain-explicit-plus}
\DeclareSIUnit{\points}{points}
\usepackage{booktabs}
\usepackage{caption}
\usepackage[numbers,sort&compress]{natbib}
\usepackage{hyperref}
\hypersetup{colorlinks=true, linkcolor=black, citecolor=black, urlcolor=blue}

\usepackage{authblk}

% Allow hyphenation inside \texttt{}. Without this, snake_case identifiers
% (mutation_diversity, consensus_divergence, ...) are unbreakable and push
% lines into the margin. [htt] enables breaking in typewriter text.
\usepackage[htt]{hyphenat}

\graphicspath{{figures/pdf/}}

% --- Shared color palette ---
% Edit these once; use everywhere via \color{cleoBlue} etc.
\definecolor{cleoBlue}{HTML}{2E5C8A}
\definecolor{cleoOrange}{HTML}{D97A29}
\definecolor{cleoGreen}{HTML}{4A8C5A}
\definecolor{cleoRed}{HTML}{B23A3A}
\definecolor{cleoGray}{HTML}{6B7280}

% --- Convenience macros ---
\newcommand{\cleo}{\textsc{cleo}}
\newcommand{\todo}[1]{\textcolor{cleoRed}{\textbf{[TODO: #1]}}}

% --- Draft / outline mode ---
% While the paper is a scaffold, the planning content is the document. Set
% \draftmodefalse to hide every \plan block and \todo marker in one shot,
% leaving only real prose — do that when preparing a submission draft.
\newif\ifdraftmode
\draftmodetrue

% Internal links (table of contents, cross-references) render in blue while
% drafting so the contents page reads as clickable. linktoc=all makes both the
% entry text and its page number active, not just the number. Submission mode
% keeps them black — see the \hypersetup on the hyperref line above.
\ifdraftmode
  \hypersetup{linktoc=all, linkcolor=cleoBlue}
\fi

% \plan{...} renders a planning note. In draft mode it is visible and set in
% gray; otherwise it vanishes.
\newcommand{\plan}[1]{%
  \ifdraftmode
    \par\vspace{0.3em}%
    \begingroup\small\color{cleoGray}\setlength{\leftskip}{1.2em}%
    #1\par\endgroup\vspace{0.3em}%
  \fi}

% planlist: bulleted planning notes. Same visibility rules as \plan.
\usepackage{comment}
\newenvironment{planlistbody}%
  {\begingroup\small\color{cleoGray}%
   \begin{itemize}\setlength{\itemsep}{0pt}\setlength{\parskip}{0pt}}%
  {\end{itemize}\endgroup}
\ifdraftmode
  \newenvironment{planlist}{\begin{planlistbody}}{\end{planlistbody}}
\else
  \excludecomment{planlist}
\fi

% --- Status tags for the document index ---
\newcommand{\statusdone}{\textcolor{cleoGreen}{\textbf{drafted}}}
\newcommand{\statusscaffold}{\textcolor{cleoOrange}{\textbf{scaffold}}}
\newcommand{\statusblocked}{\textcolor{cleoRed}{\textbf{needs data}}}

% Placeholder for an experiment result not yet measured. Visually distinct from
% \todo so the run-log table reads at a glance.
\newcommand{\todoresult}{\textcolor{cleoGray}{\textit{pending}}}


% Each section IS one panel, so a float that drifts to the next page separates
% a figure from the reading written about it. [H] pins them in place.
\usepackage{float}

\newcommand{\provenance}[1]{%
  {\footnotesize\color{cleoGray}\texttt{#1}}\par\vspace{0.4em}}

% "What it shows" against "what it can't say", side by side, because the second
% is the part that gets lost when a panel is lifted into a paper.
\newenvironment{reading}%
  {\par\vspace{0.6em}\noindent\begin{minipage}[t]{0.47\textwidth}%
   \footnotesize\textbf{What it shows.}\ }%
  {\end{minipage}\par\vspace{0.9em}}
\newcommand{\limit}[1]{%
  \end{minipage}\hfill\begin{minipage}[t]{0.47\textwidth}%
  \footnotesize\textcolor{cleoRed}{\textbf{What it can't say.}}\ #1}

\title{\cleo{}: ten campaigns, one trade}
\author{}
\date{Narrative deck --- compiled \today}

\begin{document}

\maketitle
\thispagestyle{empty}

\begin{abstract}
\noindent
One panel per experimental campaign, on matched axes where the axes are
genuinely matched. The recurring quantity is a Pareto trade: oracle pass rate
against sequence diversity, first over single designs and then over the
recombined library that the combinatorial filter actually produces. The
closing section names three narrative arcs these panels will carry, with the
panel that argues against each. This document is for choosing the argument,
not for making it.
\end{abstract}

\vspace{1em}
\noindent\textbf{How to read every panel.} Up is a library with fewer dead
wells. Right is a library covering more sequence space. $U_{15}$ is the count
of distinct (position, residue) substitutions rarefied to 15 passing designs
over 300 draws, measured against the per-position modal residue of the
LigandMPNN temperature sweep on the same backbone --- the reference under which
the published baseline comparator reproduces. Library diversity is the same
quantity over passing library members, rarefied to a matched count, so an arm
cannot win it merely by having more of its library fold.

\clearpage
\tableofcontents
\clearpage

\section{C1 --- Temperature is the dial everyone already has}
\provenance{experiments/ame/tempsweep · M0097\_1ctt · LigandMPNN, 96 designs per T}

The comparator campaign. Sampling temperature is how diversity is bought out of
LigandMPNN today, and it charges pass rate for every point of coverage.

\begin{figure}[H]\centering
  \includegraphics[width=0.82\textwidth]{figures/pdf/panel_c1.pdf}
  \caption{Six independent operating points, shaded light to dark by
  temperature. They are deliberately not joined: a line between them would
  assert a continuity nobody measured. $T = 1.0$ is omitted --- at
  \SI{4.2}{\percent} pass, too few designs survive to rarefy to 15; its mean
  pairwise Hamming is 126.3.}
\end{figure}

\begin{reading}
The best diversity the baseline reaches while still folding is
$U_{15} = \num{607}$ at \SI{26.0}{\percent} pass, at $T = 0.7$. Past that the
curve falls off a cliff.
\limit{One backbone. $T = 0.3$ scores \emph{below} $T = 0.2$ on pass rate, so
even the baseline curve is not monotone at $n = 96$ --- read the shape, not the
individual points.}
\end{reading}

\section{C2 --- The KL anchor is the dial RL adds}
\provenance{experiments/ame/klsweep\_table.py · run\_M0097\_1ctt\_cond9\_14\_klw* · 200 steps, 640 designs per arm}

Six arms of one config differing only in \texttt{kl\_weight}. The anchor holds
the policy near the pretrained model, so raising it trades pass rate back for
coverage --- the same trade as temperature, on a different mechanism.

\begin{figure}[H]\centering
  \includegraphics[width=\textwidth]{figures/pdf/panel_c2.pdf}
  \caption{Left: dashed segments show the ordering of \texttt{kl\_weight} only,
  not an interpolation. Right: raw $U$ peaks at $w = 0.02$ and then
  \emph{falls} to \num{1731} at $w = 0.10$, because so few designs pass that
  there is nothing left to count, while rarefied $U_{15}$ keeps climbing. The
  gap between those two curves is the whole argument for rarefaction. Every
  cell regenerates from \texttt{klsweep\_table.py}, the single definition of
  these numbers.}
\end{figure}

\begin{reading}
$w = 0.02$ buys $U_{15} = \num{803}$ at \SI{56.7}{\percent} pass. The dial is
real and it is smooth over four orders of magnitude of weight.
\limit{One seed per arm. Two unseeded runs of one config on this benchmark gave
\SI{68.8}{\percent} and \SI{89.6}{\percent} control pass --- a 21-point swing,
wider than the gap between the first two arms here.}
\end{reading}

\section{C3 --- The two dials are not on the same frontier}
\provenance{C1 × C2 overlaid · same backbone, same reference, same $U_{15}$ estimator}

C1 and C2 share a backbone, a reference sequence and a rarefaction depth, so
their points go on one axis pair without rescaling. This is the comparison the
paper is built on.

\begin{figure}[H]\centering
  \includegraphics[width=0.86\textwidth]{figures/pdf/panel_c3.pdf}
  \caption{The shaded quadrant is everything that dominates the baseline's best
  folding point: $U_{15} > \num{607}$ \emph{and} pass $>$ \SI{26.0}{\percent}.
  Three KL arms sit inside it; no temperature setting does.}
\end{figure}

\begin{reading}
A frontier \emph{shift}, not a point win. RL does not merely pick a better spot
on the existing trade --- it moves the achievable set outward. At $w = 0.02$ the
policy carries $1.32\times$ the baseline's coverage at $2.2\times$ its pass
rate.
\limit{Mean pairwise Hamming at $w = 0.02$ is 114.4 against the baseline's
100.5 at $T = 0.7$: comparable, not better. The win is in \emph{distinct
substitutions at matched design count}, and the claim should be stated that
narrowly.}
\end{reading}

\clearpage
\section{C4 --- Which part of the method did the work}
\provenance{experiments/ame/\ldots\_conv200\{\_legacy, \_surr, (none), \_kl\} · one backbone, four objectives}

An A/B ladder over the objective itself: the legacy reward, then the surrogate
fix, the log-prob fix, and the KL anchor on top. Each row is the row above it
plus one change.

\begin{figure}[H]\centering
  \includegraphics[width=\textwidth]{figures/pdf/panel_c4.pdf}
  \caption{Two charts rather than two $y$-axes on one. The category order is
  the order the fixes were applied, which is itself information.}
\end{figure}

\begin{reading}
The first three steps move pass rate around and barely touch coverage
($\num{468} \to \num{571} \to \num{506}$). \textbf{Only the KL anchor moves
$U_{15}$}, by \num{297} points. The diversity result is attributable to one
component.
\limit{The surrogate-fix arm is the \emph{best} on pass rate
(\SI{88.8}{\percent}) and is not the arm the paper headlines. All are single
draws; the six replicates behind this family run 571, 510, 544 with the fix and
506, 511, 410 without --- means 542 and 476, every one under the baseline's
\num{607}.}
\end{reading}

\section{C5 --- Recombining across backbones, easy to nearly impossible}
\provenance{experiments/ame/r1panel/recomb · 9 backbones · 256 recombinants per cell, AF3 best-of-5}

The first campaign at library scale: cut passing designs into fragments,
recombine across parents, fold the chimeras. Backbones are ordered by how dead
the RFdiffusion2 benchmark finds them, which is the axis along which the result
matters.

\begin{figure}[H]\centering
  \includegraphics[width=0.82\textwidth]{figures/pdf/panel_c5.pdf}
  \caption{Each line is one backbone, sampled at $T = 1.0$ to the RL policy.
  Mean \SI{53.5}{\percent} against \SI{30.5}{\percent}, RL ahead on 8 of 9. The
  single loss, in red, is M0050\_1dbt at \SI{97}{\percent} dead, where both arms
  are in single digits.}
\end{figure}

\begin{reading}
The gain is largest exactly where the benchmark is hardest --- M0255\_1mg5 at
\SI{77}{\percent} dead goes \num{25.8} to \num{79.3}. An improvement on an easy
backbone is a convenience; this is a different kind of claim.
\limit{Four backbones have no diversity reading --- too few passing members on
one arm to rarefy against --- so the diversity comparison is 6 of 8, not 9 of 9.
The two arms also differ in parent count, which is why C6 rarefies.}
\end{reading}

\clearpage
\section{C6 --- The combinatorial filter, as a Pareto}
\provenance{r1panel + fragbsweep/fragb16 ($k = 4$) + fragk ($k = 16$) · 3 backbones · all arms batch 16}

Library pass rate against library diversity, for four ways of producing the
parents. Both axes are properties of the \emph{recombined} library. This is the
panel the thesis stands on.

\begin{figure}[H]\centering
  \includegraphics[width=\textwidth]{figures/pdf/panel_c6.pdf}
  \caption{Faceted by backbone rather than overlaid: the arms cross each other
  in diversity, so one axis draws four long segments over one another and the
  shape of any single backbone becomes unreadable. Axes are shared. On
  M0315\_1ey3 the two chimera arms very nearly coincide
  (\num{1087.5}/\SI{97.3}{\percent} and \num{1083.5}/\SI{98.4}{\percent}) and
  plot as one point.}
\end{figure}

\begin{reading}
RL moves up and right of sampled on all three backbones --- a genuine
dominance. The chimera objective then moves up again (\SI{93.4}{\percent}
against \SI{76.7}{\percent}, 3 of 3) \textbf{and back to the left}: diversity
\num{895.5} against \num{1139.4}, losing on 3 of 3. Cutting at $k = 16$ instead
of $4$ recovers part of it (\num{951.6}) at the same pass rate.
\limit{No arm here dominates every other. The chimera objective trades coverage
for fidelity and \textbf{$k$ is the dial along that trade}. The paper currently
argues for the chimera arm on pass rate alone; this panel is why that has to be
stated as a position on a frontier rather than as a win.}
\end{reading}

\section{C7 --- How much of C6 could be a lucky seed}
\provenance{experiments/ame/\{r1rep,fragb32\} · T2 library eval · 256 parents + 256 recombinants per cell}

Two plain-RL replicates per backbone, differing only in \texttt{seed}, bound
what a good draw buys. The chimera arm is then placed against that band.

\begin{figure}[H]\centering
  \includegraphics[width=0.78\textwidth]{figures/pdf/panel_c7.pdf}
  \caption{Bars span the two plain-RL seeds; the point is the chimera arm.
  Margins over the top of the band are $+25.8$, $+52.0$ and $+34.8$ points ---
  three to six times the largest seed-to-seed difference.}
\end{figure}

\begin{reading}
This is why a single chimera seed is worth something here when single-seed
results usually are not: the effect clears the whole measured band on every
backbone.
\limit{The band is measured on plain-RL/RF3 arms. \textbf{No pool arm has ever
been run twice}, so the chimera point's own error bar is unmeasured. The
contrast varies objective \emph{and} oracle together and cannot apportion the
gain. Quote the recombinant pass rate, never retention --- retention swings
$0.767 \to 1.611$ on an unchanged config.}
\end{reading}

\clearpage
\section{C8 --- Whether the training number was ever honest}
\provenance{experiments/ame/\{fragbsweep,fragb16,fragb32\}/offline\_probe · 38 held-out AF3 probe points}

The reward sits inside a structure predictor, so a training pass rate is not
evidence on its own. Gap $=$ oracle-balanced training level minus held-out AF3
on chimeras rebuilt deterministically from the same step's parents.

\begin{figure}[H]\centering
  \includegraphics[width=\textwidth]{figures/pdf/panel_c8.pdf}
  \caption{Above the zero line the training number overstates. Pooled over all
  38 points: mean $+0.015$, median $-0.000$, with 7 of 38 overstating by more
  than $0.10$. The pooled mean is carried by the two smallest cells, so it is
  quoted with its breakdown or not at all.}
\end{figure}

\begin{reading}
At batch 16 and 32 the in-loop number is honest to within $0.005$. The one
contrast that isolates the oracle --- \texttt{fragb32} against
\texttt{fragb32rf3}, both batch 32, both chimera-objective, both self-resumed
--- puts \textbf{the oracle pool at $0.046$} of avoided overstatement.
\limit{The hatched arm, \texttt{fragrl4} at $+0.279$, warm-starts from a
\texttt{fragrl3} checkpoint while the pool arms start fresh. It varies
initialisation as well as oracle: it is the arm the project started from, never
the controlled contrast. Neither arm in the controlled pair is replicated in
seed.}
\end{reading}

\section{C9 --- The only campaign with a wet lab in it}
\provenance{adaptyv\_library · de novo PETase, 209\,aa, 5 fragments · 3,135 wells assayed}

128 protein parts over five slots spanning $1.1 \times 10^{7}$ recombinants;
1,920 BsaI DNA blocks, one-pot assembly, one well per recombinant.
\textbf{\SI{67}{\percent} registered above background at a median of 29
mutations from the parent.}

\begin{figure}[H]\centering
  \includegraphics[width=\textwidth]{figures/pdf/panel_c9.pdf}
  \caption{Left: scoring the right state at the right atoms is worth roughly
  $2\times$ the global number. But distance from the parent ranks activity at
  $-0.65$, better than any AF3 metric, so the zero-shot claim is made only on
  the \num{2539} designs $\geq 20$ mutations out, where active-site pLDDT still
  gives $+0.59$. Right: median activity by distance bin, log scale.}
\end{figure}

\begin{reading}
The independence assumption that \texttt{results/fragments.tex} names as
load-bearing survives a physical test. Searching the library does move activity:
median $2.2 \to 12.2$ across four rounds, best recombinant $20.5 \to 31.7$, with
the highest of all four ($47.8$) in round~2.
\limit{No wild-type comparator was assayed, so no fold-over-parent number
exists --- the archived $250\times$ was always a placeholder. The 20 designs
scoring above 100 are a different assay run with no shared normalising control,
so nothing crosses that line. And on this scaffold recombination at depth bought
reachable diversity and a \SI{67}{\percent} register rate, \emph{not} a better
enzyme than local work.}
\end{reading}

\clearpage
\section{C10 --- The same pipeline on a second enzyme, and it inverts}
\provenance{paper/figures/data/petase\_super\_*\_transfer.csv · second de novo hydrolase, 160\,aa, 2,509 purified}

Each point is one metric, plotted by how well it predicted activity on PETase
against how well it predicted on \texttt{super}. On the diagonal the metric
transfers; in the off-diagonal quadrants it points the wrong way.

\begin{figure}[H]\centering
  \includegraphics[width=0.86\textwidth]{figures/pdf/panel_c10.pdf}
  \caption{AF3 concepts: $\rho$-of-$\rho$s $= -0.604$ over 49 shared concepts,
  with only 10 of 49 agreeing in sign against ${\sim}24$ expected by chance.
  All ten of PETase's strongest predictors invert, mean $+0.470 \to -0.100$.
  Rosetta fails more mildly --- $+0.035$, 33 of 66 --- i.e. no transferable
  signal rather than an inverted one.}
\end{figure}

\begin{reading}
One family survives: PETase's three strongest Rosetta predictors are all
His--Asp distance terms, and all three \textbf{keep their sign} while collapsing
7--21$\times$ in magnitude. Catalytic-geometry strain points the same way on
both enzymes and nearly stops mattering --- the only concept family that does.
\limit{This is the wet-lab counterpart of C8, and it is the strongest argument
in the paper for the library strategy: a predictor's agreement with the quantity
you care about is a property of the specific system, re-earned per target, never
inherited. Because zero-shot scores do not transfer, selection has to be done by
experiment --- which is what a combinatorially compatible library at a
\SI{67}{\percent} register rate is for.}
\end{reading}

\clearpage
\section{What a library costs}
\provenance{design calculation, not a campaign · Twist eBlocks at \$0.07 per base · 200\,aa enzyme}

Every campaign above is eventually answerable to this one. A protein of $L$
amino acids is $3L$ bases. Cut it into $k$ equal fragments and each fragment is
$3L/k$ bases, so a budget $B$ at a price $c$ per base buys

\begin{equation*}
  v \;=\; \frac{B}{3Lc} \quad\text{variants per slot,}
  \qquad
  S \;=\; v^{k} \;=\; \left(\frac{B}{3Lc}\right)^{k} \quad\text{recombinants.}
\end{equation*}

\noindent
There is no $k$ in $v$, and that is the whole point: cutting into $k$ pieces
buys $k$ times as many fragments, each $1/k$ as long, and the two cancel
exactly. \textbf{$k$ is a pure exponent on a number set entirely by budget.}
Doubling the budget multiplies the library by $2^{k}$ --- at $k = 6$, twice the
money is 64 times the library.

For a \num{200}-aa enzyme at \$0.07 per base, one full-length
equivalent costs \$42, so \$\num{10000} stocks \num{238} variants in every
slot regardless of how finely it is cut.

\begin{figure}[H]\centering
  \includegraphics[width=0.92\textwidth]{figures/pdf/panel_cost_size.pdf}
  \caption{Both axes are logarithmic, so a different price per base slides
  every curve sideways without changing its shape --- the figure is reusable at
  any quote. The PETase library is shown twice because it was ordered twice
  over: \num{128} distinct protein parts realised as \num{1920} DNA parts, a
  $15\times$ redundancy that is the dominant term in what was actually spent.}
\end{figure}

\begin{reading}
Synthesis is not the binding constraint, and it is not close. At \$\num{10000}
--- less than the folding compute behind a single campaign above --- the
affordable library at $k = 6$ is $\num{1.8e14}$ recombinants, five orders of
magnitude past what one-pot assembly can transform and eleven past what the
PETase screen actually read. The budget buys sequence space nobody can look at.
\limit{Three things bind first, all of them visible in the earlier panels.
\textbf{Parents:} $v$ variants per slot needs $v$ passing designs, and past
\$\num{10752} the \num{256}-parent libraries of C6 and C7 run out before the
money does. \textbf{Assembly:} ${\sim}10^{9}$ transformants. \textbf{Screening:}
\num{3135} wells, in the one screen this project has run. Cost sets the ceiling;
it has never set the floor.}
\end{reading}

\begin{figure}[H]\centering
  \includegraphics[width=\textwidth]{figures/pdf/panel_cost_k.pdf}
  \caption{Left: the idealised exponent against what vendor terms actually
  deliver. Right: a fragment is billed at $\max(m,\, 3L/k + h)$, for an order
  minimum $m$ and $h$ bases of Golden Gate sites and overhangs. On a
  \num{200}-aa enzyme the minimum binds from $k = 3$ onward, and by $k = 8$
  three quarters of every base bought is padding. \textbf{$m$ and $h$ are
  assumed} --- \num{300} and \num{40} bases --- and are the two numbers to
  check against a live quote before anyone plans a spend from this figure. The
  \$0.07 per base is not assumed.}
\end{figure}

\begin{reading}
The order minimum is what converts $k$ from free to expensive. Idealised, $k$
costs nothing; after the minimum, going from $k = 4$ to $k = 16$ at
\$\num{10000} drops the variants per slot from \num{119} to \num{30}, and the
library still grows --- $\num{2e8}$ to $\num{3.8e23}$ --- because the exponent
beats the division. The penalty is real and it loses anyway.
\limit{So this calculation does \emph{not} pick a $k$. Within everything
synthesis charges for, larger $k$ always wins, which means $k$ has to be chosen
on the axis C6 measures --- where fidelity and coverage trade --- and on how
many junctions an assembly will actually tolerate. The cost model only says
that money will not make that choice for you.}
\end{reading}

\clearpage
\section{Three arcs these panels will carry}

Each arc names the panels that carry it and the panel that argues against it.

\paragraph{Arc 1 --- RL moves the frontier.}
C1 $\to$ C2 $\to$ C3 as the spine, C5 for breadth, C6 for the library payoff.
\emph{Against it:} C4, where only one of four components moves coverage at all,
and the best pass rate belongs to an arm the paper does not headline.

\paragraph{Arc 2 --- fragments are the unit.}
C6 $\to$ C7 as the spine, C9 for physical realizability. \emph{Against it:} C6
again --- the chimera objective loses diversity on 3 of 3 backbones, so
``recombinable'' is bought, not free.

\paragraph{Arc 3 --- you cannot predict, so you must order.}
C8 $\to$ C10 as the spine, with C6 as the thing you order and C9 as proof it
assembles. \emph{Against it:} nothing in these panels. It rests on the two
strongest results in the project --- 38 probe points and two independent enzyme
screens --- and it makes the diversity/fidelity frontier the deliverable rather
than the claim.

\vspace{1em}
\noindent\textbf{What is not regenerable.} Table~1's \cleo{} rows cannot be
rebuilt: \texttt{backbone19} and \texttt{selection2} were pruned and no archive
survives, which also breaks \texttt{ame\_diversity.py},
\texttt{ame\_frontier.py} and \texttt{ame\_novelty.py}. The panels here avoid
those trees entirely.

\end{document}
